%---------------------------------------------------
% کلاس سند: article برای اسناد متنی، با اندازه کاغذ A4 و فونت پایه 16
\documentclass[a4paper,16pt]{article}

% بسته‌های ریاضی برای معادلات و نمادهای پیشرفته
\usepackage{amsmath, amssymb, amstext}

% بسته رنگ برای استفاده از رنگ در متن یا محیط‌ها
\usepackage{xcolor}

% بسته گرافیک برای درج تصاویر
\usepackage{graphicx}

% تنظیمات حاشیه صفحه
\usepackage{geometry}

% بسته لینک‌دهی (هایپرلینک‌ها)
\usepackage{hyperref}

% بسته صفحه‌آرایی برای تنظیم هدر و فوتر
\usepackage{fancyhdr}


% بسته xepersian برای پشتیبانی از زبان فارسی با XeLaTeX
\usepackage{xepersian}

% تعیین فونت فارسی (فایل فونت باید موجود باشد یا در سیستم نصب باشد)
\settextfont{B Nazanin}

% تنظیم اندازه حاشیه‌ها از هر طرف برابر با 1 اینچ
\geometry{margin=1in}

% تنظیم سبک fancy برای سربرگ و پاورقی
\pagestyle{fancy}

% سمت چپ هدر: مشخص‌کننده سری تمرین و نام درس
\lhead{تمرین سری چهار درس مکانیک تحلیلی }

% وسط هدر: نام دانشکده
\chead{دانشکده فیزیک دانشگاه صنعتی خواجه نصیرالدین طوسی}

% سمت راست هدر: تاریخ روز به‌صورت خودکار
\rhead{\today}

% پایین صفحه: شماره صفحه
\cfoot{\thepage}

% شروع سند
\begin{document}
	
	% صفحه عنوان (بدون هدر و فوتر)
	\thispagestyle{empty}
	\begin{center}
		% لوگوی دانشگاه (باید فایل تصویر در مسیر پروژه باشد)
			\includegraphics[width=0.5\textwidth]{../../../images/image-E_1/image-1.png} \\[10pt]
		
		% عنوان سند
		\textbf{\LARGE عنوان : تمرین سری چهار}\\[20pt]
		
		% نیم‌سال تحصیلی
		\textbf{\LARGE نیم‌سال تحصیلی : 4041}\\[10pt]
		
		% نام مدرس درس
		\textbf{\Large مدرس : دکتر محمد انصاری فرد}\\[10pt]
		
		% مبحث تمرین 
		\textbf{\Large مبحث تمرین : نوسانگر ها }\\[10pt]
		
		% مهلت تحویل تمرین
		\textbf{\Large مهلت تحویل : 21 آذر}
	\end{center}
	
	% صفحه جدید برای ادامه سند
	\newpage
	
	% ایجاد فهرست مطالب به‌صورت خودکار (براساس \section و \subsection و ...)
	\tableofcontents
	\newpage
	
	% شروع بخش سوالات - سوال اول
\section{سوال اول}
جابجایی $x(t)$ و سرعت $v(t)$ (یا $\dot{x}(t)$) نوسانگر بیش‌میرا \lr{overdamped} را به صورت توابع هذلولوی بیان کنید.


\section{سوال دوم}
نشان دهید که برای نوسانگر با میرایی کم \lr{(lightly damped)}، \(Q \approx \dfrac{\omega_0}{\Delta\omega}\).


\section{سوال سوم}
یک نوسانگر خطی میرایافته \lr{damped linear oscillator}، که ابتدا در موقعیت تعادل خود ساکن است، تحت تأثیر تابع تحریک زیر قرار می‌گیرد:

\[
\frac{F(t)}{m} =
\begin{cases}
	0, & t<0\\[2mm]
	a \frac{t}{\tau}, & 0<t<\tau\\[1mm]
	a, & t>\tau
\end{cases}
\]

پاسخ سیستم \lr{response function} را بیابید. حال \(\tau \to 0\) را در نظر بگیرید و نشان دهید که جواب به حالت تابع پله \lr{step function} تبدیل می‌شود.


\section{سوال چهارم}
یک خودرو با جرم \(1000\) \lr{kg} (به‌همراه سرنشینان) طوری است که به ازای اضافه شدن هر \(100\) \lr{kg} جرم، به اندازه \(1.0\) \lr{cm} بیشتر نسبت به سطح جاده نشست می‌کند.  
خودرو با مولفه افقی سرعت ثابت \(20\) \lr{km/h} روی جاده‌ای حرکت می‌کند که به‌صورت یک موج سینوسی با برجستگی‌هایی به دامنه \(5.0\) \lr{cm} و طول موج \(20\) \lr{cm} مدل شده است. فاصله میان چرخ جلو و عقب \(2.4\) \lr{m} است.

فرض کنید حرکت عمودی خودرو را می‌توان مانند یک نوسانگر هارمونیکِ تحتِ نیروی محرک و بدون میرایی مدل کرد. جرم چرخ‌ها و فنرها را نادیده بگیرید و فرض کنید چرخ‌ها همیشه با سطح جاده در تماس‌اند.  
دامنه نوسان عمودی خودرو را به‌دست آورید.


\section{سوال پنجم}
دو جرم با مقدار \lr{g} \(m_1 = 100\)\,و \lr{g} \(m_2 = 200\)\,روی یک مسیر افقی بدون اصطکاک به‌صورت آزاد حرکت می‌کنند و توسط یک فنر با ثابت نیروی \lr{N/m}\(k = 0.5\)\, به یکدیگر متصل شده‌اند. بسامد نوسان این دستگاه را بیابید.


	
	\vspace{10pt}
\textbf{موفق باشید.}
	% پایان سند
\end{document}
